\chapterimage{figs/chapterhead.png}
\chapter{Núcleo do simulador}
\label{chap:nucleo_do_simulador}

\section{Introdução}
\label{sec:cap7_introducao}

Se o capítulo anterior organizou o mapa geral do projeto, este capítulo entra na região mais importante desse mapa: o núcleo do simulador. É aqui que o \textit{GenESyS} deixa de ser visto apenas como aplicação e passa a ser entendido como sistema que representa modelos, agenda eventos, controla repetições, mantém o tempo simulado e coordena a observabilidade do experimento.

O estudo do kernel é decisivo por uma razão simples: tudo o que o usuário faz na GUI e tudo o que o desenvolvedor adiciona por parser, \textit{plugins} ou aplicações depende, em última instância, da infraestrutura definida aqui. Não importa se o modelo foi aberto visualmente, descrito textualmente ou construído por código. Em todos os casos, ele precisará ser representado, verificado, executado e observado por classes centrais do kernel.

Neste capítulo, o foco recairá sobre um conjunto de classes e estruturas particularmente importantes:
\begin{itemize}
    \item \texttt{Simulator};
    \item \texttt{Model};
    \item \texttt{ModelSimulation};
    \item \texttt{OnEventManager};
    \item lista de eventos futuros;
    \item entidades e eventos correntes;
    \item mecanismos de controle do ciclo da simulação.
\end{itemize}

O objetivo não é ainda esgotar todas as classes do kernel, mas construir uma compreensão firme da dinâmica central do GenESyS.

\section{O kernel como infraestrutura de simulação}
\label{sec:cap7_kernel_como_infraestrutura}

O kernel do \textit{GenESyS} deve ser entendido como a camada que torna possível a existência do simulador independentemente da interface usada. A GUI, a aplicação terminal, modelos textuais, \textit{plugins} e ferramentas podem variar; o kernel, porém, é a base que sustenta a representação do modelo e a execução da simulação.

Em termos arquiteturais, isso significa que o kernel precisa resolver pelo menos quatro problemas fundamentais:

\begin{enumerate}
    \item como representar um modelo de simulação;
    \item como representar o estado da execução desse modelo;
    \item como avançar no tempo simulado por meio de eventos;
    \item como tornar a simulação observável e controlável.
\end{enumerate}

Esses quatro problemas aparecem com clareza nas classes centrais do projeto. É por isso que estudar o kernel não é estudar ``mais um pacote'' do sistema, mas a parte que define o comportamento essencial do simulador.

\section{A classe \texttt{Simulator} como ponto de entrada do kernel}
\label{sec:cap7_classe_simulator}

A classe \texttt{Simulator} ocupa posição especial na visão do kernel. Ela se apresenta como a classe principal do sistema de simulação e fornece acesso a gerenciadores importantes como \texttt{PluginManager}, \texttt{ModelManager}, \texttt{TraceManager}, \texttt{ParserManager} e \texttt{ExperimentManager}. Essa organização já revela algo importante: o GenESyS não foi estruturado em torno de uma única classe monolítica, mas em torno de um núcleo que coordena subsistemas especializados.

Do ponto de vista do desenvolvedor, a classe \texttt{Simulator} é útil por dois motivos. Primeiro, porque oferece uma visão sintética da arquitetura do projeto. Segundo, porque funciona como ponto de acesso de alto nível aos demais mecanismos do sistema. Em vez de entrar no software por fragmentos isolados, o leitor pode começar pela percepção de que existe uma classe central que articula o ecossistema do simulador.

\subsection{O que essa classe sugere sobre a arquitetura}
\label{subsec:cap7_o_que_simulator_sugere}

A mera presença dos gerenciadores acessíveis a partir de \texttt{Simulator} já indica cinco pilares arquiteturais:
\begin{itemize}
    \item o sistema administra modelos;
    \item o sistema administra \textit{plugins};
    \item o sistema possui uma infraestrutura explícita de tracing;
    \item o sistema possui um subsistema de parser;
    \item o sistema prevê suporte a experimentos.
\end{itemize}

Mesmo antes de ler os detalhes de cada gerenciador, o desenvolvedor já percebe que o núcleo do GenESyS foi pensado como infraestrutura coordenadora e não como bloco único de execução.

\section{A classe \texttt{Model} como centro do modelo de simulação}
\label{sec:cap7_classe_model}

A classe \texttt{Model} merece atenção especial porque ela representa, no kernel, o próprio modelo de simulação. Sua importância é explicitada já na documentação do cabeçalho: ela é descrita como provavelmente a classe mais importante do kernel do GenESyS. Essa formulação faz sentido. É nela que se articulam componentes, \textit{data definitions}, lista de eventos futuros, controles, respostas, persistência, parser, verificação de consistência, gerenciamento de entidades e acesso ao estado da simulação.

Em termos práticos, \texttt{Model} é a classe na qual o modelo deixa de ser apenas ideia ou diagrama e passa a ser objeto manipulável pelo software.

\subsection{Por que \texttt{Model} é tão central}
\label{subsec:cap7_por_que_model_central}

A centralidade de \texttt{Model} decorre do fato de que ela articula vários mecanismos essenciais ao mesmo tempo:
\begin{itemize}
    \item inserção e remoção de elementos do modelo;
    \item criação e remoção de entidades;
    \item envio de entidades entre componentes;
    \item verificação de consistência do modelo;
    \item persistência e carregamento;
    \item acesso a parser, tracing e simulação;
    \item manutenção da lista de eventos futuros.
\end{itemize}

Isso significa que a classe \texttt{Model} não é apenas um contêiner de componentes. Ela é o centro organizador da vida do modelo dentro do simulador.

\subsection{Modelo como combinação de componentes e dados}
\label{subsec:cap7_modelo_componentes_dados}

A própria documentação da classe indica que o modelo é representado principalmente por:
\begin{itemize}
    \item uma coleção de componentes de modelo, adequadamente configurados e conectados;
    \item uma coleção de \textit{model data definitions}, que sustentam a infraestrutura do modelo.
\end{itemize}

Essa distinção é arquiteturalmente fundamental. Ela antecipa um princípio que voltará em capítulos posteriores: o modelo não é feito apenas de blocos de fluxo. Ele depende também de uma camada de dados e definições que dá suporte à simulação.

\subsection{A lista de eventos futuros}
\label{subsec:cap7_lista_eventos_futuros}

Entre os atributos de \texttt{Model}, um merece destaque especial: a lista de eventos futuros. A própria documentação do código a descreve como um dos elementos mais importantes de um sistema de simulação orientado a eventos. Isso se justifica facilmente. É a partir dessa lista que a simulação avança cronologicamente, tomando sempre o próximo evento a ser processado.

Para o desenvolvedor, isso é essencial. O tempo simulado não avança como simples laço incremental arbitrário. Ele avança por disparo ordenado de eventos agendados, e a lista de eventos futuros é o mecanismo que sustenta esse processo.

\section{A classe \texttt{ModelSimulation}}
\label{sec:cap7_classe_modelsimulation}

Se \texttt{Model} representa o modelo, \texttt{ModelSimulation} representa a execução desse modelo. Essa distinção é extremamente importante. O modelo e sua simulação são objetos relacionados, mas não idênticos. O primeiro descreve a estrutura do sistema modelado; o segundo controla seu ciclo de execução.

A documentação de \texttt{ModelSimulation} é particularmente esclarecedora. Ela a descreve como a camada madura e funcional de execução da simulação, baseada em repetições, comprimento de replicação, configuração de \textit{warm-up} e processamento de eventos. Também enfatiza que, embora existam abstrações mais altas de experimentação em desenvolvimento, \texttt{ModelSimulation} continua sendo o núcleo estável de execução. Essa observação é extremamente valiosa para o desenvolvedor, porque indica qual classe realmente concentra hoje o comportamento operacional maduro da simulação.

\subsection{Responsabilidades de \texttt{ModelSimulation}}
\label{subsec:cap7_responsabilidades_modelsimulation}

As responsabilidades de \texttt{ModelSimulation} podem ser organizadas em alguns grupos principais:

\begin{itemize}
    \item iniciar, pausar, avançar passo a passo e parar a simulação;
    \item controlar repetições;
    \item manter o tempo simulado corrente;
    \item administrar comprimento de replicação e período de \textit{warm-up};
    \item verificar condições de término;
    \item gerenciar breakpoints;
    \item disparar relatórios;
    \item coordenar o processamento de eventos.
\end{itemize}

Observe-se que essa classe não é mero detalhe operacional. Ela concentra a lógica efetiva da execução do modelo e, por isso, deve ser estudada com cuidado por qualquer desenvolvedor que pretenda alterar a dinâmica central do simulador.

\subsection{O ciclo da simulação}
\label{subsec:cap7_ciclo_da_simulacao}

A leitura da classe \texttt{ModelSimulation} permite perceber uma estrutura lógica clara do ciclo de simulação:

\begin{enumerate}
    \item inicialização da simulação;
    \item inicialização de uma replicação;
    \item processamento sequencial de eventos;
    \item verificação de pontos de interrupção e condições de término;
    \item encerramento de replicação;
    \item encerramento da simulação.
\end{enumerate}

Esse ciclo é muito importante para a compreensão do sistema. Ele mostra que a execução não é uma simples chamada monolítica de ``run'', mas um processo estruturado em fases, cada uma com responsabilidades próprias.

\subsection{Tempo simulado, replicações e \textit{warm-up}}
\label{subsec:cap7_tempo_replications_warmup}

Um dos méritos da classe \texttt{ModelSimulation} é tornar explícito que a execução do modelo envolve mais do que a simples passagem de eventos. Há também a gestão de:
\begin{itemize}
    \item tempo simulado corrente;
    \item número de repetições;
    \item comprimento de replicação;
    \item período de \textit{warm-up};
    \item unidade de tempo de base;
    \item condição de término.
\end{itemize}

Esses elementos revelam que a classe foi desenhada para sustentar simulação orientada a experimentação e análise, e não apenas execução ad hoc de um fluxo de eventos.

\section{Eventos, evento corrente e processamento sequencial}
\label{sec:cap7_eventos_evento_corrente_processamento}

A simulação discreta orientada a eventos depende de uma ideia central: o sistema evolui de evento em evento. No \textit{GenESyS}, essa lógica aparece de forma muito clara na relação entre \texttt{Model}, \texttt{ModelSimulation}, \texttt{Event} e a lista de eventos futuros.

Durante a simulação, há sempre um evento corrente em processamento. Esse evento define o instante simulado atual e o contexto da operação em curso. Ao ser processado, ele pode produzir efeitos sobre entidades, componentes, dados do modelo e a própria agenda futura do sistema.

Essa organização tem duas implicações importantes para o desenvolvedor:
\begin{enumerate}
    \item o estado do sistema não é atualizado de forma difusa, mas em pontos discretos de processamento;
    \item a semântica dos componentes e de grande parte das ações do modelo depende do evento corrente.
\end{enumerate}

Essa observação se conecta diretamente ao parser, a entidades, a componentes e a várias expressões ligadas ao estado do modelo, que serão discutidas em capítulos posteriores.

\section{A classe \texttt{OnEventManager} e a observabilidade do núcleo}
\label{sec:cap7_classe_oneventmanager}

Se \texttt{ModelSimulation} controla a execução, \texttt{OnEventManager} torna essa execução observável por meio de eventos de alto nível. Esse é um ponto arquitetural muito elegante do GenESyS e merece destaque especial.

A classe \texttt{OnEventManager} permite registrar tratadores para diversos tipos de ocorrência do ciclo de vida do modelo e da simulação. Entre eles estão:
\begin{itemize}
    \item sucesso na verificação do modelo;
    \item carregamento e salvamento de modelo;
    \item início e término de replicação;
    \item processamento de evento;
    \item criação, movimentação e remoção de entidades;
    \item início, pausa, retomada e fim da simulação;
    \item disparo de breakpoint.
\end{itemize}

Essa infraestrutura mostra que o GenESyS trata observabilidade como parte constitutiva do kernel, e não como mero recurso periférico de interface.

\subsection{Por que isso é arquiteturalmente importante}
\label{subsec:cap7_por_que_oneventmanager_importante}

O \texttt{OnEventManager} separa elegantemente duas coisas:
\begin{itemize}
    \item a lógica de execução da simulação;
    \item a lógica de reação a eventos relevantes da simulação.
\end{itemize}

Essa separação é valiosa porque permite que aplicações, ferramentas, relatórios e mecanismos de depuração se conectem ao comportamento do kernel sem precisar misturar sua lógica diretamente ao processamento central. Trata-se, em essência, de um mecanismo de observação e reação baseado em eventos do próprio simulador.

\subsection{ModelEvent e SimulationEvent}
\label{subsec:cap7_modelevent_simulationevent}

A presença das classes auxiliares \texttt{ModelEvent} e \texttt{SimulationEvent} reforça essa arquitetura. Elas funcionam como objetos de contexto para os tratadores registrados, carregando informações sobre o modelo, o tempo simulado, o evento corrente, a replicação, o estado de pausa ou parada, a entidade criada e o componente de destino, entre outros dados relevantes.

Com isso, o \texttt{OnEventManager} não apenas anuncia que ``algo ocorreu'', mas fornece contexto suficiente para que a reação a esse evento seja tecnicamente útil.

\section{Breakpoints e controle fino da execução}
\label{sec:cap7_breakpoints_controle_fino_execucao}

A classe \texttt{ModelSimulation} também torna explícita a existência de listas de breakpoints por tempo, componente e entidade. Essa característica é particularmente relevante porque mostra que o kernel foi concebido não apenas para executar modelos até o fim, mas para permitir observação controlada da simulação.

Do ponto de vista do usuário, isso aparece na GUI como capacidade de pausar ou avançar passo a passo. Do ponto de vista do desenvolvedor, isso revela que o controle fino da execução não é adicionado pela interface de forma superficial. Ele está previsto na própria infraestrutura do kernel.

Essa é uma decisão arquitetural importante. Ela significa que depuração, inspeção e experimentação controlada não dependem exclusivamente da aplicação gráfica; elas têm fundamento no núcleo de simulação.

\section{Uma leitura integrada do núcleo}
\label{sec:cap7_leitura_integrada_nucleo}

A melhor forma de compreender o núcleo do \textit{GenESyS} é integrar as funções das classes estudadas até aqui.

\begin{itemize}
    \item \texttt{Simulator} coordena o acesso a subsistemas centrais.
    \item \texttt{Model} representa o modelo, seus componentes, dados e eventos futuros.
    \item \texttt{ModelSimulation} controla o ciclo de execução desse modelo.
    \item \texttt{OnEventManager} torna os eventos relevantes observáveis e reagíveis.
\end{itemize}

Essa integração ajuda a perceber que o kernel não é um conjunto de classes paralelas sem articulação. Ele é uma infraestrutura coerente na qual cada classe ocupa papel específico dentro da mesma dinâmica de simulação.

\begin{table}[htbp]
    \centering
    \caption{Papéis das classes centrais do núcleo do GenESyS.}
    \label{tab:cap7_papeis_classes_centrais}
    \small
    \begin{tabular}{|p{3.6cm}|p{7.6cm}|}
        \hline
        \textbf{Classe} & \textbf{Papel principal} \\
        \hline
        \texttt{Simulator} & Classe de mais alto nível do kernel, coordenando acesso a gerenciadores centrais \\
        \hline
        \texttt{Model} & Representação do modelo de simulação, de seus componentes, dados, entidades e eventos futuros \\
        \hline
        \texttt{ModelSimulation} & Controle da execução do modelo, incluindo repetições, tempo simulado, condições de término e processamento de eventos \\
        \hline
        \texttt{OnEventManager} & Registro e disparo de tratadores associados a eventos relevantes do modelo e da simulação \\
        \hline
    \end{tabular}
\end{table}

\section{Considerações Finais}
\label{sec:cap7_consideracoes_finais}

Este capítulo apresentou o núcleo do simulador a partir de suas classes mais centrais. O ponto principal a reter é que o \textit{GenESyS} organiza sua infraestrutura de simulação de forma clara: há uma classe de alto nível que coordena subsistemas, uma classe que representa o modelo, uma classe que controla a execução desse modelo e um mecanismo explícito de observabilidade orientada a eventos. Essa organização é suficientemente robusta para sustentar aplicações, parser, \textit{plugins}, tracing, breakpoints e futuras extensões.

Com essa base estabelecida, o próximo passo natural é avançar para a estrutura dos próprios elementos do modelo. Em outras palavras, se este capítulo mostrou como o sistema representa e executa um modelo, o capítulo seguinte mostrará de que tipos de objetos esse modelo é feito: componentes, \textit{data definitions} e mecanismos de extensibilidade por \textit{plugins}.

Teste seu entendimento desse conteúdo respondendo as seguintes perguntas:

\begin{enumerate}
    \item Qual é a diferença conceitual entre \texttt{Model} e \texttt{ModelSimulation} no kernel do GenESyS?

    \item Por que a lista de eventos futuros é uma estrutura tão importante na simulação orientada a eventos?

    \item Que papel o \texttt{OnEventManager} desempenha na observabilidade do sistema?

    \item Por que a presença de breakpoints no kernel é uma decisão arquitetural relevante e não apenas um detalhe de interface?
\end{enumerate}